\chapter{Referencja metadata}

Ten rozdzial zbiera pola metadata wspierane przez MoonChunk i podpowiada, kiedy ich
uzywac. Nie wszystkie beda potrzebne w kazdym projekcie. Warto jednak wiedziec, ze
jezyk wspiera szeroki zakres konfiguracji dokumentu.

\begin{longtable}{>{\raggedright\arraybackslash}p{0.25\textwidth}>{\raggedright\arraybackslash}p{0.65\textwidth}}
\toprule
\textbf{Pole} & \textbf{Znaczenie i typowe zastosowanie} \\
\midrule
\endfirsthead
\toprule
\textbf{Pole} & \textbf{Znaczenie i typowe zastosowanie} \\
\midrule
\endhead
\bottomrule
\endfoot
lang & Jezyk dokumentu, np. \mc{"pl"} albo \mc{"en"}. Wplywa m.in. na interpretacje tresci przez przegladarke i narzedzia wspomagajace. \\

dir & Kierunek tekstu, zwykle \mc{"ltr"} dla jezykow pisanych od lewej do prawej. \\

htmlClass & Klasy CSS przypisywane do znacznika \mc{<html>}. Przydatne do globalnego sterowania wygladem. \\

charset & Kodowanie znakow, zwykle \mc{"utf-8"}. W praktyce wartosc niemal obowiazkowa. \\

viewport & Ustawienia widoku dla urzadzen mobilnych, np. \mc{"width=device-width, initial-scale=1"}. \\

title & Tytul dokumentu pokazywany na karcie przegladarki i wykorzystywany przez wyszukiwarki. \\

metaDescription & Krótki opis strony dla SEO i podgladu w wynikach wyszukiwania. \\

metaKeywords & Lista slow kluczowych. Wspolczesnie ma mniejsze znaczenie niz dawniej, ale nadal moze porzadkowac projekt. \\

metaAuthor & Autor strony albo treści. \\

metaRobots & Wskazowki dla robotow indeksujacych, np. \mc{"index,follow"}. \\

themeColor & Kolor motywu wykorzystywany przez niektore przegladarki i urzadzenia mobilne. \\

canonicalUrl & Kanoniczny adres URL strony. Przydatny, gdy jedna tresc wystepuje pod wieloma adresami. \\

faviconHref & Sciezka do favikony strony. \\

appleTouchIconHref & Ikona dla urzadzen Apple. \\

manifestHref & Sciezka do manifestu aplikacji webowej. \\

ogType & Typ Open Graph, np. \mc{"website"} albo \mc{"article"}. \\

ogTitle & Tytul widoczny przy udostepnianiu strony w social media. \\

ogDescription & Opis strony dla Open Graph. \\

ogImage & Obraz do podgladu przy udostepnianiu. \\

ogUrl & Kanoniczny adres strony dla Open Graph. \\

ogSiteName & Nazwa serwisu. \\

ogLocale & Lokalizacja Open Graph, np. \mc{"pl\_PL"}. \\

twitterCard & Typ karty Twitter/X, np. \mc{"summary"} albo \mc{"summary\_large\_image"}. \\

twitterSite & Konto serwisu w mediach spolecznosciowych. \\

twitterCreator & Konto autora tresci. \\

twitterTitle & Tytul do karty Twitter/X. \\

twitterDescription & Opis do karty Twitter/X. \\

twitterImage & Obraz do karty Twitter/X. \\

preloadLinks & Wlasne znaczniki \mc{<link rel='preload' ...>} wstawiane do \mc{<head>}. \\

preconnectLinks & Wlasne znaczniki \mc{<link rel='preconnect' ...>}. \\

styles & Linki do CSS lub inne style, zwykle jako gotowy fragment HTML. \\

headScripts & Skrypty umieszczane w sekcji \mc{<head>}. \\

headExtra & Dodatkowy surowy HTML dla \mc{<head>}. \\

bodyClass & Klasy CSS dla znacznika \mc{<body>}. \\

pageId & Identyfikator strony trafiajacy do atrybutu danych w ciele dokumentu. \\

topBar & Fragment HTML wstawiany nad glowna trescia. \\

header & Fragment HTML naglowka strony. \\

footer & Fragment HTML stopki strony. \\

modals & Miejsce na kontenery modali i innych elementow nakladanych. \\

scripts & Skrypty umieszczane pod koniec dokumentu, zwykle w \mc{<body>}. \\

bodyEndExtra & Dodatkowy HTML wstawiany na koncu \mc{<body>}. \\
\end{longtable}

\section{Praktyczna rada}
Nie ma potrzeby uzupelniac wszystkich pol od razu. Zacznij od kilku najwazniejszych:
\mc{lang}, \mc{charset}, \mc{viewport}, \mc{title}, \mc{metaDescription}, \mc{styles}.
Potem rozbudowuj dokument o kolejne informacje w miare potrzeb projektu.
